\RequirePackage{plautopatch}
\documentclass[platex, a4paper, 11pt]{jsarticle}
\usepackage{jpConstract}

\title{団体規約　例}
\date{}


\begin{document}
\maketitle

% ====================
\shou{総則}
% ====================
\jou{名称}
\jou{所在地}
\jou{目的}
  \kou{もくてき}
  \kou{もくてき}
    \gou{goal}
\jou{活動}
  \kou[]{数字だけ隠れてほしい}
\jou{組織}
  \kou{本会は、前条の活動を行うために、次の部門を置く（ここの数字は隠れてほしい）}
    \gou{実行委員会}
    \gou{首相会}
    \gou{企画委員会}
  \kou{この項と第1項で上下から号を挟みたい}
    \gou{無理だった。}
    \gou{号の下に項が出力されてしまう}
  \kou{第2項以降は問題なさそう}
    \gou{問題なし}
\jou{項が隠れてほしい}
  \kou[first]{テスト}

% \section{これは非常に長いセクションタイトルで、ページの右端を超えても自動的に折り返されることを確認します。LaTeXのデフォルトの設定であれば、特に何も変更しなくても適切に動作します。}

% このセクションでは、タイトルの折り返しについて説明します。

% \rulechapter{総則}
%   \rulearticle{名称}
%   \rulearticle{所在地}
%   \rulearticle{目的}
%   \rulearticle{活動}
%     \ruleparagraph{前条の目的を達成するために，下記の行事を開催する．}
%       \ruleitem{関西学生基礎スキー大会}
%       \ruleitem{実行委員会が企画した行事}
%       \ruleitem{企画委員が企画した行事}
%       \ruleitem{その他本会の目的を達成するために必要な事項}
%     \ruleparagraph{これは非常に長いセクションタイトルで、ページの右端を超えても自動的に折り返されることを確認します。LaTeXのデフォルトの設定であれば、特に何も変更しなくても適切に動作します。}
%   \rulearticle{組織}
%   \ruleparagraph{本会は，前条の活動を行うために，次の部門を置く．}
%     \ruleitem{実行委員会}
%     \ruleitem{主将会}
%     \ruleitem{企画委員会}
% \rulechapter{参加クラブ}
%   \rulearticle{名称}
%   \rulearticle{所在地}
%   \rulearticle{目的}
%   \rulearticle{活動}
%   \rulearticle{組織}

% \begin{enumerate}
%   \item aaaaaaaaaaaaaaaaaaaaaaaaaaaaaaaaaaaaaaaaaaaaaaaaaaaaaaaaaaaaaaaaaaaaaaaaaaaaaaaaaaaaaaaaaaaaaaaaaaaaaaaaaaaaaaaaaaaaaaaaaa
%   \item aaaaaa
% \end{enumerate}

% \begin{Fusoku}
%   \fusoku{asf}
% \end{Fusoku}

\end{document}